\documentclass[full-course-notes.tex]{subfiles}

\begin{document}
\lecturemeta{3}
  {Tokens, Patterns, and Lexemes; the Role of the Lexical Analyzer; Lexical Errors}
  {CLO 1}
  {\S3.1 (The Role of the Lexical Analyzer), \S3.3 (Specification of Tokens -- terminology preview)}
  {[AP] Modern Compiler Implementation, Ch.\,2}
\lecshort{Tokens, Patterns \& Lexemes}
\setcounter{section}{0}
\makelecturetitle

\begin{coreidea}[Lecture Roadmap]
\begin{enumerate}
  \item Where lexical analysis sits, and exactly how it talks to the parser.
  \item Why lexer and parser are two separate phases instead of one.
  \item The three core vocabulary words of this unit: \textbf{token}, \textbf{pattern},
  \textbf{lexeme} --- with several fully worked examples.
  \item Token \textbf{attributes} and what actually gets passed on to the parser.
  \item \textbf{Lexical errors}: what a lexer really can and cannot catch, and how it recovers.
\end{enumerate}
This lecture builds the vocabulary and mental model. \emph{How} patterns are written --
regular expressions and finite automata -- is Lecture~\ref{lec:4} onward.
\end{coreidea}

\section{Where Lexical Analysis Fits}
\label{sec:l3-where-lexical-fits}

\dbtag The \textbf{lexical analyzer} (also called a \textbf{scanner} or \textbf{lexer}) is the
first phase of a compiler. It reads the raw stream of characters making up the source program
and groups them into a stream of \textbf{tokens} --- meaningful chunks like keywords,
identifiers, numbers, and operators --- that the rest of the compiler works with instead of raw
text.

\begin{center}
\begin{tikzpicture}[
  box/.style={draw, thick, rounded corners, minimum height=1.1cm, minimum width=3cm, align=center, font=\small},
  tool/.style={box, fill=dbGreenBg, draw=dbGreen},
  data/.style={box, fill=nuLight, draw=nuNavy},
  >=Latex
]
  \node[data] (src) at (0,0) {Source Program\\(character stream)};
  \node[tool] (lex) at (4.2,0) {Lexical\\Analyzer};
  \node[tool] (parser) at (8.6,0) {Parser};
  \node[data] (st) at (4.2,-2.2) {Symbol Table};

  \draw[->, thick] (src) -- (lex);
  \draw[->, thick] ([yshift=2mm]lex.east) -- node[above, font=\scriptsize]{token} ([yshift=2mm]parser.west);
  \draw[->, thick] ([yshift=-2mm]parser.west) -- node[below, font=\scriptsize]{getNextToken()} ([yshift=-2mm]lex.east);
  \draw[<->, thick, dashed] (lex) -- (st);
\end{tikzpicture}
\end{center}

\begin{coreidea}[How the Lexer and Parser Actually Talk to Each Other]
In nearly every real compiler, the lexer does \emph{not} tokenize the whole file up front and
hand the parser a list. Instead, the parser drives the process: whenever the parser needs another
token, it calls a function conventionally named \texttt{getNextToken()}. The lexer resumes
scanning exactly where it left off, produces the next single token, and returns it. This
``pull'' model, sometimes called treating the lexer as a \textbf{coroutine} of the parser, avoids
ever materializing the full token list in memory and lets the lexer and parser interleave their
work one token at a time.
\end{coreidea}

\subsection{Why Not Just One Phase?}
\label{sec:l3-why-separate}

\dbtag DB gives three concrete engineering reasons for splitting lexical analysis out as its own
phase rather than folding it into the parser:

\begin{enumerate}
  \item \textbf{Simplicity of design.} Separating out low-level tasks (recognizing keywords,
  numbers, string literals) lets the parser work purely with grammatical structure, using
  tokens as its atomic symbols, instead of also worrying about individual characters.
  \item \textbf{Improved efficiency.} Specialized, table-driven techniques for recognizing
  tokens (built from finite automata, covered once we reach Lecture~\ref{lec:4} and beyond) are
  considerably faster than the general-purpose parsing techniques would be if forced to operate
  character-by-character.
  \item \textbf{Enhanced portability.} Input-device and operating-system-specific quirks ---
  end-of-line conventions (\texttt{\textbackslash n} vs.\ \texttt{\textbackslash r\textbackslash n}), character
  encodings, special end-of-file characters --- are confined entirely to the lexer. The parser,
  and everything after it, is completely insulated from these differences.
\end{enumerate}

\begin{examfocus}
This ``three reasons'' list (simplicity, efficiency, portability) is a favorite short-answer
question. Be able to name and briefly justify all three, not just one.
\end{examfocus}

\section{Tokens, Patterns, and Lexemes}
\label{sec:l3-tokens-patterns-lexemes}

\dbtag These three terms are used precisely and are easy to blur together. Getting the
distinction exactly right is one of the most exam-tested ideas in this entire unit.

\begin{defbox}{Token, Pattern, Lexeme}
\begin{itemize}
  \item A \textbf{token} is a pair \texttt{<token-name, attribute-value>} representing a
  \emph{category} of lexical unit, e.g.\ \texttt{ID}, \texttt{NUM}, \texttt{RELOP}. The
  token-name is an abstract symbol used by the parser.
  \item A \textbf{pattern} is the rule that describes the \emph{set of possible lexemes} that
  can produce a given token -- for identifiers, informally, ``a letter followed by any number of
  letters and digits.'' Lecture~\ref{lec:4} makes patterns precise using regular expressions.
  \item A \textbf{lexeme} is the actual sequence of characters in the source program that
  \emph{matches} a pattern and is classified as an instance of some token.
\end{itemize}
\end{defbox}

\begin{center}
\begin{tikzpicture}[
  box/.style={draw, thick, rounded corners, minimum height=1cm, minimum width=3.1cm, align=center, font=\small},
  >=Latex
]
  \node[box, fill=nuLight, draw=nuNavy] (chars) at (0,0) {Source\\Characters};
  \node[box, fill=enrichGoldBg, draw=enrichGold] (pat) at (4,0) {matched against\\a \textbf{Pattern}};
  \node[box, fill=dbGreenBg, draw=dbGreen] (lex) at (8,0) {yields a\\\textbf{Lexeme}};
  \node[box, fill=examRedBg, draw=examRed] (tok) at (12,0) {emits a\\\textbf{Token}\\\texttt{<name, attr>}};
  \draw[->, thick] (chars) -- (pat);
  \draw[->, thick] (pat) -- (lex);
  \draw[->, thick] (lex) -- (tok);
\end{tikzpicture}
\end{center}

\begin{examplebox}{DB-Style Token/Pattern/Lexeme Table}
\renewcommand{\arraystretch}{1.3}
\begin{tabularx}{\textwidth}{@{}l X l@{}}
\toprule
\textbf{Token} & \textbf{Informal Description of Pattern} & \textbf{Sample Lexemes} \\
\midrule
\texttt{WHILE} & the exact keyword \texttt{while} & \texttt{while} \\
\texttt{ID} & a letter, followed by any number of letters or digits & \texttt{i}, \texttt{sum}, \texttt{score2} \\
\texttt{NUM} & any numeric constant & \texttt{0}, \texttt{42}, \texttt{3.14} \\
\texttt{RELOP} & \texttt{<}, \texttt{<=}, \texttt{=}, \texttt{<>}, \texttt{>}, or \texttt{>=} & \texttt{<=}, \texttt{<>} \\
\texttt{LITERAL} & any characters between a pair of double quotes, except a double quote itself & \texttt{"Total = "} \\
\bottomrule
\end{tabularx}
\end{examplebox}

\subsection{Worked Example 1: Tokenizing a Function Call}
\label{sec:l3-worked-example-1}

Consider the following line of C-like source code:

\begin{center}
\ttfamily\large
\colorbox{dbGreenBg}{printf}\colorbox{examRedBg}{(}\colorbox{enrichGoldBg}{"Total = \%d\textbackslash n"}\colorbox{examRedBg}{,}\ \colorbox{dbGreenBg}{score}\colorbox{examRedBg}{)}\colorbox{examRedBg}{;}
\end{center}
\begin{center}
\small\begin{tabular}{@{}p{2.6cm}p{2.6cm}p{2.6cm}p{2.6cm}@{}}
\colorbox{dbGreenBg}{\phantom{x}}\ identifier &
\colorbox{examRedBg}{\phantom{x}}\ punctuation/operator &
\colorbox{enrichGoldBg}{\phantom{x}}\ string literal &
\end{tabular}
\end{center}

\begin{center}
\renewcommand{\arraystretch}{1.25}
\begin{tabularx}{\textwidth}{@{}l l X@{}}
\toprule
\textbf{Lexeme} & \textbf{Token Name} & \textbf{Attribute-Value (passed to parser)} \\
\midrule
\texttt{printf} & \texttt{ID} & pointer to symbol-table entry for \texttt{printf} \\
\texttt{(} & \texttt{LPAREN} & none needed \\
\texttt{"Total = \%d\textbackslash n"} & \texttt{LITERAL} & pointer to the string's stored value \\
\texttt{,} & \texttt{COMMA} & none needed \\
\texttt{score} & \texttt{ID} & pointer to symbol-table entry for \texttt{score} \\
\texttt{)} & \texttt{RPAREN} & none needed \\
\texttt{;} & \texttt{SEMI} & none needed \\
\bottomrule
\end{tabularx}
\end{center}

\section{Attributes for Tokens}
\label{sec:l3-attributes}

\dbtag When more than one lexeme can match a token's pattern (as with \texttt{ID} or
\texttt{NUM}), the lexer must pass along enough extra information -- the \textbf{attribute-value}
-- for later phases to know exactly \emph{which} identifier or \emph{which} number was matched.
For identifiers, this is essentially always a pointer into the symbol table.

\begin{examplebox}{Worked Example 2: An Assignment Statement}
Source line: \quad \texttt{E = M * C ** 2}

Assume \texttt{**} is this language's exponentiation operator. The lexer's output, as a stream
of \texttt{<token-name, attribute-value>} pairs, is:
\begin{center}
\ttfamily
\begin{tabular}{@{}l@{}}
\ttfamily <ID, ptr to symtab entry for E> \\
<ASSIGN\_OP,\ -- >  \\
<ID, ptr to symtab entry for M> \\
<MULT\_OP,\ -- > \\
<ID, ptr to symtab entry for C> \\
<EXP\_OP,\ -- > \\
<NUM, integer value 2>
\end{tabular}
\end{center}
Notice: \texttt{ID} tokens all carry a symbol-table pointer as their attribute so that later
phases can distinguish \texttt{E} from \texttt{M} from \texttt{C}; the operator tokens need no
attribute at all, since the token-name alone (\texttt{ASSIGN\_OP} vs.\ \texttt{MULT\_OP}) already
says everything the parser needs to know.
\end{examplebox}

\begin{examfocus}
Given a short program line, be able to produce the exact stream of
\texttt{<token-name, attribute-value>} pairs, in order, the way both worked examples above do.
This is the single most common long-answer question style for this unit.
\end{examfocus}

\subsection{Worked Example 3: A Complete Statement Block}
\label{sec:l3-worked-example-3}

\begin{center}
\begin{minipage}{0.5\textwidth}
\ttfamily
int sum = 0;\\
while (i <= 10) \{\\
\ \ sum = sum + i;\\
\ \ i = i + 1;\\
\}
\end{minipage}
\end{center}

\begin{center}
\renewcommand{\arraystretch}{1.15}
\footnotesize
\begin{tabularx}{\textwidth}{@{}l l X | l l X@{}}
\toprule
\textbf{Lexeme} & \textbf{Token} & \textbf{Attribute} & \textbf{Lexeme} & \textbf{Token} & \textbf{Attribute} \\
\midrule
\texttt{int} & \texttt{TYPE} & -- & \texttt{sum} & \texttt{ID} & ptr(\texttt{sum}) \\
\texttt{sum} & \texttt{ID} & ptr(\texttt{sum}) & \texttt{=} & \texttt{ASSIGN} & -- \\
\texttt{=} & \texttt{ASSIGN} & -- & \texttt{sum} & \texttt{ID} & ptr(\texttt{sum}) \\
\texttt{0} & \texttt{NUM} & 0 & \texttt{+} & \texttt{PLUS} & -- \\
\texttt{;} & \texttt{SEMI} & -- & \texttt{i} & \texttt{ID} & ptr(\texttt{i}) \\
\texttt{while} & \texttt{WHILE} & -- & \texttt{;} & \texttt{SEMI} & -- \\
\texttt{(} & \texttt{LPAREN} & -- & \texttt{i} & \texttt{ID} & ptr(\texttt{i}) \\
\texttt{i} & \texttt{ID} & ptr(\texttt{i}) & \texttt{=} & \texttt{ASSIGN} & -- \\
\texttt{<=} & \texttt{RELOP} & LE & \texttt{i} & \texttt{ID} & ptr(\texttt{i}) \\
\texttt{10} & \texttt{NUM} & 10 & \texttt{+} & \texttt{PLUS} & -- \\
\texttt{)} & \texttt{RPAREN} & -- & \texttt{1} & \texttt{NUM} & 1 \\
\texttt{\{} & \texttt{LBRACE} & -- & \texttt{;} & \texttt{SEMI} & -- \\
 & & & \texttt{\}} & \texttt{RBRACE} & -- \\
\bottomrule
\end{tabularx}
\end{center}

Two details worth noticing in this longer example:
\begin{itemize}
  \item \texttt{sum} and \texttt{i} each get \emph{their own} symbol-table pointer the first
  time they're seen, and the \emph{same} pointer every time they reappear -- the symbol table
  guarantees one entry per distinct identifier, no matter how many times it's used.
  \item Whitespace (the spaces and newlines between tokens) and indentation produce \emph{no}
  tokens at all -- they are consumed and discarded by the lexer, along with any comments, before
  the parser ever sees anything.
\end{itemize}

\begin{coreidea}[The Lexer's ``Housekeeping'' Duties]
Beyond producing tokens, DB notes the lexer is also the natural place to: strip out whitespace
and comments (neither produces a token); and track line numbers (and often column numbers) so
that later phases can report errors with accurate source locations, e.g.\ ``\texttt{undeclared
variable `x' at line 42}.''
\end{coreidea}

\section{Lexical Errors}
\label{sec:l3-lexical-errors}

\dbtag It is surprisingly rare for a lexer to be able to detect an error entirely on its own,
\emph{by itself, in isolation} -- because almost any short character sequence is a \emph{prefix}
of some valid lexeme for some token.

\begin{examplebox}{The Classic Illustration}
Consider the fragment \texttt{fi (a == f(x))} in a C-like language, where the programmer
\emph{meant} to type the keyword \texttt{if}. The lexer cannot know this. \texttt{fi} is a
perfectly legal lexeme for token \texttt{ID} (``a letter followed by letters/digits''), so the
lexer happily emits \texttt{<ID, ptr(fi)>}. The mistake will only be caught later -- by the
\emph{parser}, if \texttt{fi} appears somewhere no identifier is grammatically valid, or by
\emph{semantic analysis}, if \texttt{fi} is used as if it were a function without ever having
been declared.
\end{examplebox}

Genuine lexical errors are ones where \emph{no pattern for any token} can match the remaining
input at all, or where a token's own internal structure is malformed:

\begin{itemize}
  \item An illegal character that starts no valid token in the language, e.g.\ a stray
  \texttt{\#} or \texttt{@} in a language that assigns those characters no meaning.
  \item An unterminated string literal: opening quote seen, but the line (or file) ends before a
  matching closing quote.
  \item A malformed numeric literal, e.g.\ \texttt{3.14.15} (two decimal points), or
  \texttt{2r} immediately followed by a letter with no valid meaning in that position.
  \item An unterminated comment: a \texttt{/*} seen with no matching \texttt{*/} anywhere before
  end of file.
\end{itemize}

\subsection{Error Recovery Strategies}
\label{sec:l3-error-recovery}

\dbtag A production-quality lexer does not simply stop at the first error --- it tries to
recover and keep going, so the programmer can see \emph{several} distinct errors from one
compilation attempt rather than fixing them one at a time.

\begin{center}
\begin{tikzpicture}[
  box/.style={draw, thick, rounded corners, minimum height=0.9cm, minimum width=3.6cm, align=center, font=\small},
  dec/.style={draw, thick, diamond, aspect=2.4, minimum width=3.4cm, minimum height=1.4cm, align=center, font=\small, fill=nuLight, draw=nuNavy},
  >=Latex
]
  \node[box, fill=nuLight, draw=nuNavy] (start) at (0,0) {Read next\\characters};
  \node[dec] (match) at (0,-2.1) {Matches some\\token's pattern?};
  \node[box, fill=dbGreenBg, draw=dbGreen] (emit) at (5,-2.1) {Emit token,\\advance input};
  \node[box, fill=examRedBg, draw=examRed] (err) at (0,-4.6) {Report lexical\\error};
  \node[box, fill=enrichGoldBg, draw=enrichGold] (panic) at (0,-6.6) {Panic-mode: delete\\char(s) until a valid\\token can resume};

  \draw[->, thick] (start) -- (match);
  \draw[->, thick] (match) -- node[above, font=\scriptsize]{yes} (emit);
  \draw[->, thick] (match) -- node[left, font=\scriptsize]{no} (err);
  \draw[->, thick] (err) -- (panic);
  \draw[->, thick] (emit.north) to[out=90,in=0] (start.east);
  \draw[->, thick] (panic.west) to[out=180,in=270] (start.south);
\end{tikzpicture}
\end{center}

\dbtag DB lists several concrete local-correction strategies a recovering lexer can attempt,
beyond plain panic mode:

\begin{itemize}
  \item \textbf{Deleting} an extraneous character.
  \item \textbf{Inserting} a missing character.
  \item \textbf{Replacing} one incorrect character with a correct one.
  \item \textbf{Transposing} two adjacent characters.
\end{itemize}

\begin{historynote}
DB points toward a more principled version of this idea: choosing whichever single correction
requires the \emph{fewest} character edits -- insertions, deletions, substitutions -- to turn
the bad input into something valid. This is exactly \textbf{edit distance} (Levenshtein
distance), the same string-similarity measure used today in spell-checkers and DNA sequence
alignment. Aho and Peterson described a minimum-distance error-correcting parser as early as
1972; the idea that ``the best fix is the cheapest fix'' turns out to generalize far beyond
lexical analysis, into syntax error recovery and even runtime auto-correct systems.
\end{historynote}

\begin{examfocus}
Be able to explain, with a concrete example like \texttt{fi (a == f(x))}, \emph{why} a
misspelled keyword is usually \textbf{not} a lexical error even though it looks like an obvious
typo. Then be able to list at least two genuine, purely lexical errors (unterminated string,
illegal character) and describe panic-mode recovery in one sentence.
\end{examfocus}

\section{Beyond the Dragon Book: How Real Lexers Handle Harder Cases}
\label{sec:l3-beyond-db}

\begin{enrichment}[The Reserved-Word Trick]
A naive lexer might need a separate hand-written check for every keyword (\texttt{if},
\texttt{while}, \texttt{return}, \dots) before falling back to the generic identifier pattern.
Production lexers instead use a much simpler trick: pre-load the symbol table with every
reserved word \emph{before} scanning begins, each marked with its own keyword token. The lexer
then has only \emph{one} pattern for anything ``letter followed by letters/digits.'' When it
matches such a lexeme, it looks the lexeme up in the symbol table: if an entry already exists
(because it's a reserved word), the lexer returns \emph{that} token (e.g.\ \texttt{WHILE}); if
no entry exists, it's a genuine user identifier, so the lexer creates a new symbol-table entry
and returns \texttt{ID}. One pattern, one lookup, no special-casing in the automaton itself.
\end{enrichment}

\begin{enrichment}[Maximal Munch -- A Preview]
When several patterns could all match a prefix of the remaining input, real lexers apply the
\textbf{longest-match} (``maximal munch'') rule: always take the \emph{longest} lexeme that
matches \emph{any} pattern. This is why \texttt{<=} is scanned as a single \texttt{RELOP} token
and not as \texttt{<} followed by \texttt{=}, and why \texttt{i} in \texttt{int} is not
prematurely cut off as an identifier before the lexer has looked far enough ahead to see the
rest of the word. C's \texttt{-{}-{}} (decrement) versus \texttt{-} \texttt{-} (two unary
minuses) and \texttt{-{}>} (member access through a pointer) are further classic examples where
getting the lookahead right matters. We formalize \emph{why} this rule is needed, and how a
scanner implements it efficiently, once we cover finite automata in Lecture~\ref{lec:4} and
beyond.
\end{enrichment}

\begin{enrichment}[Numeric Literals Are Far Messier Than ``a Sequence of Digits'']
DB's \texttt{NUM} pattern is deliberately simplified. Real languages support numeric literal
formats far beyond plain decimal integers, and each variant is a distinct lexical pattern the
scanner must recognize:
\begin{center}
\small
\begin{tabularx}{\textwidth}{@{}l X@{}}
\toprule
\textbf{Format} & \textbf{Example} \\
\midrule
Hexadecimal / octal / binary & \texttt{0x1F}, \texttt{0o17}, \texttt{0b1010} \\
Scientific notation & \texttt{6.022e23}, \texttt{1.5E-10} \\
Digit-group separators & \texttt{1\_000\_000} (Python, Java, Rust, C++14) \\
Type suffixes & \texttt{3.14f} (float), \texttt{100L} (long), \texttt{42u} (unsigned) \\
Complex/imaginary literals & \texttt{3+4j} (Python) \\
\bottomrule
\end{tabularx}
\end{center}
Each of these is its own regular pattern (or family of patterns) that a real scanner's
specification must include -- one more reason lexer-generator tools (covered once we reach the
lexer-generator lecture later in the course) are preferred over hand-written code once a
language's literal syntax grows this rich.
\end{enrichment}

\begin{enrichment}[String Literals: Escapes, Raw Strings, and Interpolation]
String-literal lexing hides real complexity:
\begin{itemize}
  \item \textbf{Escape sequences} (\texttt{\textbackslash n}, \texttt{\textbackslash t},
  \texttt{\textbackslash\textbackslash}, \texttt{\textbackslash uXXXX}) must be recognized
  \emph{inside} the literal's pattern, not treated as ending it.
  \item \textbf{Raw strings} (Python's \texttt{r"..."}, C\#'s \texttt{@"..."}) deliberately
  \emph{disable} escape processing -- a single string-literal token category can therefore need
  two different sub-patterns.
  \item \textbf{String interpolation / f-strings} (Python's \texttt{f"Total = \{score\}"},
  JavaScript template literals) are the hardest case: the lexer must recognize where an embedded
  \emph{expression} begins inside a string and, in effect, briefly hand control back to a
  full expression lexer/parser before resuming the string literal. This breaks the assumption
  that lexical analysis is a single flat pass over characters, and is a genuinely active area of
  parser/lexer-design discussion in modern language implementations.
\end{itemize}
\end{enrichment}

\begin{enrichment}[Not Every Language Is Purely Character-Driven: Indentation as Tokens]
Python and Haskell use indentation, not braces, to delimit blocks. Their lexers cannot rely on
simple per-character regular patterns alone -- they maintain an explicit \textbf{indentation
stack} and synthesize special \texttt{INDENT} and \texttt{DEDENT} pseudo-tokens whenever a
logical line's indentation increases or decreases relative to the stack's top, in addition to
the ordinary tokens for the line's actual content. This is a clean example of \emph{why} DB's
model of ``patterns matched against a flat character stream'' is the right \emph{starting}
point but not the end of the story for every real language.
\end{enrichment}

\begin{enrichment}[Unicode Identifiers]
DB's ``a letter followed by letters and digits'' quietly assumes ASCII. Modern language
standards (Python 3, Swift, Julia, Rust as of recent editions) permit \emph{any} Unicode
character in the ``Letter'' category as part of an identifier -- so accented Latin letters
(\texttt{caf\'e}), non-Latin scripts (Greek, Arabic, Devanagari, Han characters), and even emoji
are all valid identifiers in the languages that permit them (Swift famously allows an emoji as a
variable name). Real lexer specifications reference the Unicode standard's character-category
tables rather than a simple \texttt{[A-Za-z]} range.
\end{enrichment}

\begin{enrichment}[Seeing Real Tokenizers at Work]
You do not have to take any of this on faith -- most language toolchains expose their tokenizer
directly:
\begin{itemize}
  \item Python: \texttt{python3 -m tokenize myfile.py} prints every token, its exact source
  position, and its type.
  \item Clang: \texttt{clang -Xclang -dump-tokens -fsyntax-only file.c} dumps the full token
  stream Clang's own lexer produced.
  \item Node.js/Acorn or Babel's own \texttt{--tokens} flags do the same for JavaScript.
\end{itemize}
Running one of these on a file with a deliberate typo -- an unterminated string, an unexpected
character -- is the fastest way to see exactly what a real lexer reports as an error versus what
it silently accepts and passes on to the parser.
\end{enrichment}

\section{Self-Check Questions}
\label{sec:l3-selfcheck}

\begin{enrichment}[Practice -- Not Graded]
\begin{enumerate}
  \item For the statement \texttt{total = total + rate * 60;}, write out the full stream of
  \texttt{<token-name, attribute-value>} pairs, following the style of Worked Examples 2 and 3.
  \item Explain why \texttt{2x} (digit immediately followed by a letter) can be treated as a
  genuine lexical error in many languages, while \texttt{fi} cannot, even though both look like
  ``obvious mistakes'' to a human reader.
  \item A lexer encounters an opening \texttt{"} with no matching closing quote before the end
  of the file. Describe two different, reasonable ways a lexer could recover from this and keep
  scanning the rest of the file.
  \item Why must the longest-match (maximal munch) rule apply when scanning \texttt{i <= 10}, so
  that \texttt{<=} is not mistakenly split into \texttt{<} followed by \texttt{=}?
  \item Explain, in your own words, why Python's lexer cannot be built from character-level
  patterns alone the way a C lexer can.
\end{enumerate}
\end{enrichment}

\begin{takeaways}
\begin{itemize}
  \item A \textbf{token} is an abstract \texttt{<name, attribute>} pair; a \textbf{pattern} is
  the rule describing which lexemes belong to a token; a \textbf{lexeme} is the actual matched
  source text. Keep these three cleanly separated.
  \item Lexer and parser are split into separate phases for \textbf{simplicity, efficiency, and
  portability} -- know all three reasons.
  \item The lexer and parser typically interact via a pull-based \texttt{getNextToken()} call,
  not by materializing a full token list up front.
  \item Token attributes -- almost always a symbol-table pointer for \texttt{ID} -- are how
  later phases recover exactly \emph{which} identifier or literal a token represents.
  \item Most apparent ``typos'' (like \texttt{fi} for \texttt{if}) are \emph{not} lexical errors,
  because they are valid lexemes for some other token; genuine lexical errors are things no
  pattern can match at all (illegal characters, unterminated strings/comments, malformed
  numbers).
  \item Real-world lexers go well beyond DB's simplified patterns: the reserved-word trick,
  maximal munch, rich numeric/string literal formats, indentation-based tokens, and Unicode
  identifiers are all standard in production compilers.
\end{itemize}
\end{takeaways}

\section*{References for This Lecture}
\begin{thebibliography}{99}
\bibitem[DB]{db} Aho, A.~V., Lam, M.~S., Sethi, R., and Ullman, J.~D. \textit{Compilers:
Principles, Techniques, and Tools}. 2nd~ed. Pearson, 2006.
\bibitem[AP]{ap} Appel, A.~W. \textit{Modern Compiler Implementation in C/Java/ML}. Cambridge
University Press, 2004.
\end{thebibliography}

\end{document}
